\documentclass[12pt,letterpaper]{article}
\usepackage{graphicx,textcomp}
\usepackage{natbib}
\usepackage{setspace}
\usepackage{fullpage}
\usepackage{color}
\usepackage[reqno]{amsmath}
\usepackage{amsthm}
\usepackage{fancyvrb}
\usepackage{amssymb,enumerate}
\usepackage[all]{xy}
\usepackage{endnotes}
\usepackage{lscape}
\newtheorem{com}{Comment}
\usepackage{float}
\usepackage{hyperref}
\newtheorem{lem} {Lemma}
\newtheorem{prop}{Proposition}
\newtheorem{thm}{Theorem}
\newtheorem{defn}{Definition}
\newtheorem{cor}{Corollary}
\newtheorem{obs}{Observation}
\usepackage[compact]{titlesec}
\usepackage{dcolumn}
\usepackage{tikz}
\usetikzlibrary{arrows}
\usepackage{multirow}
\usepackage{xcolor}
\newcolumntype{.}{D{.}{.}{-1}}
\newcolumntype{d}[1]{D{.}{.}{#1}}
\definecolor{light-gray}{gray}{0.65}
\usepackage{url}
\usepackage{listings}
\usepackage{color}

\definecolor{codegreen}{rgb}{0,0.6,0}
\definecolor{codegray}{rgb}{0.5,0.5,0.5}
\definecolor{codepurple}{rgb}{0.58,0,0.82}
\definecolor{backcolour}{rgb}{0.95,0.95,0.92}

\lstdefinestyle{mystyle}{
	backgroundcolor=\color{backcolour},   
	commentstyle=\color{codegreen},
	keywordstyle=\color{magenta},
	numberstyle=\tiny\color{codegray},
	stringstyle=\color{codepurple},
	basicstyle=\footnotesize,
	breakatwhitespace=false,         
	breaklines=true,                 
	captionpos=b,                    
	keepspaces=true,                 
	numbers=left,                    
	numbersep=5pt,                  
	showspaces=false,                
	showstringspaces=false,
	showtabs=false,                  
	tabsize=2
}
\lstset{style=mystyle}
\newcommand{\Sref}[1]{Section~\ref{#1}}
\newtheorem{hyp}{Hypothesis}


\title{Problem Set 4}
\date{Due: November 26, 2026}
\author{Applied Stats/Quant Methods 1}


\begin{document}
	\maketitle
	\section*{Instructions}
	\begin{itemize}
		\item Please show your work! You may lose points by simply writing in the answer. If the problem requires you to execute commands in \texttt{R}, please include the code you used to get your answers. Please also include the \texttt{.R} file that contains your code. If you are not sure if work needs to be shown for a particular problem, please ask.
		\item Your homework should be submitted electronically on GitHub.
		\item This problem set is due before 23:59 on Thursday November 26, 2026. No late assignments will be accepted.
	\end{itemize}



	\vspace{.5cm}
\section*{Question 1: Economics}
\vspace{.25cm}
\noindent 	
In this question, use the \texttt{prestige} dataset in the \texttt{car} library. First, run the following commands:

\begin{verbatim}
install.packages(car)
library(car)
data(Prestige)
help(Prestige)
\end{verbatim} 


\noindent We would like to study whether individuals with higher levels of income have more prestigious jobs. Moreover, we would like to study whether professionals have more prestigious jobs than blue and white collar workers.

\newpage
\begin{enumerate}
	
	\item [(a)]
	Create a new variable \texttt{professional} by recoding the variable \texttt{type} so that professionals are coded as $1$, and blue and white collar workers are coded as $0$ (Hint: \texttt{ifelse}).
	
	\begin{doublespace}

	{\footnotesize
		\begin{verbatim}
	education          income          women           prestige         census       type   
	Min.   : 6.380   Min.   :  611   Min.   : 0.000   Min.   :14.80   Min.   :1113   bc  :44  
	1st Qu.: 8.445   1st Qu.: 4106   1st Qu.: 3.592   1st Qu.:35.23   1st Qu.:3120   prof:31  
	Median :10.540   Median : 5930   Median :13.600   Median :43.60   Median :5135   wc  :23  
	Mean   :10.738   Mean   : 6798   Mean   :28.979   Mean   :46.83   Mean   :5402   NAs : 4  
	3rd Qu.:12.648   3rd Qu.: 8187   3rd Qu.:52.203   3rd Qu.:59.27   3rd Qu.:8312            
	Max.   :15.970   Max.   :25879   Max.   :97.510   Max.   :87.20   Max.   :9517 
	\end{verbatim}}
	\lstinputlisting[language=R, firstline=39, lastline=41]{PS04_answersJX.R}

	\end{doublespace}
	
	\newpage
	
	\item [(b)]
	Run a linear model with \texttt{prestige} as an outcome and \texttt{income}, \texttt{professional}, and the interaction of the two as predictors (Note: this is a continuous $\times$ dummy interaction.)
	
	\begin{doublespace}
	\lstinputlisting[language=R, firstline=44, lastline=45]{PS04_answersJX.R}
	
	{\footnotesize
	\begin{verbatim}
		Call:
		lm(formula = prestige ~ income + professionals + income:professionals, 
		data = Prestige)
		
		Residuals:
		Min      1Q  Median      3Q     Max 
		-14.852  -5.332  -1.272   4.658  29.932 
		
		Coefficients:
		Estimate Std. Error t value Pr(>|t|)    
		(Intercept)          21.1422589  2.8044261   7.539 2.93e-11 ***
		income                0.0031709  0.0004993   6.351 7.55e-09 ***
		professionals        37.7812800  4.2482744   8.893 4.14e-14 ***
		income:professionals -0.0023257  0.0005675  -4.098 8.83e-05 ***
		---
		Signif. codes:  0 ‘***’ 0.001 ‘**’ 0.01 ‘*’ 0.05 ‘.’ 0.1 ‘ ’ 1
		
		Residual standard error: 8.012 on 94 degrees of freedom
		(4 observations deleted due to missingness)
		Multiple R-squared:  0.7872,	Adjusted R-squared:  0.7804 
		F-statistic: 115.9 on 3 and 94 DF,  p-value: < 2.2e-16
	\end{verbatim}
	}
	\end{doublespace}
	
	\newpage
	\item [(c)]
	Write the prediction equation based on the result.
	
	\begin{align*}
		\hat{y} &= \hat{\beta}_0 + \hat{\beta}_1 X_1 + \hat{\beta}_2 D_i + \hat{\beta}_3 X_1 D_i \\
	\end{align*}
	
	General prediction equation.
	\begin{align*}
		\hat{prestige} &= 21.1423 + 0.003171income + 37.7813professional - 0.002326income \times professional
	\end{align*}
	
	Prediction equation for professionals (D = 1):
	\begin{align*}
		\hat{prestige} &= 58.9235 + 0.000845income
	\end{align*}
	
	Prediction equation for white/blue collar(non-professionals, D = 0) :
	\begin{align*}
		\hat{prestige} &= 21.1423 + 0.003171income
	\end{align*}


	
\newpage
	\item [(d)]
	Interpret the coefficient for \texttt{income}.\\
	Answer: \\
	\begin{doublespace}
		For non-professional jobs, one unit increase in income is associated with 0.003171 unit increase in prestige score, on average.
	\end{doublespace}
	
	\item [(e)]
	Interpret the coefficient for \texttt{professional}.\\	
	Answer: \\
	\begin{doublespace}
		When income is 0, professional jobs are predicted to have 37.7813 unit in prestige score higher than non-professional jobs, on average.
	\end{doublespace}
	
	\newpage
	\item [(f)]
	What is the effect of a \$1,000 increase in income on prestige score for professional occupations? In other words, we are interested in the marginal effect of income when the variable \texttt{professional} takes the value of $1$. Calculate the change in $\hat{y}$ associated with a \$1,000 increase in income based on your answer for (c).\\
	
	Answer: \\
	Prediction equation for professionals (D = 1):
	\begin{align*}
		\hat{prestige_{D=1}} &= 21.1423 + 0.003171income + 37.7813 - 0.002326income \\
		\hat{prestige_{D=1}} &= 58.9235 + 0.000845income \\
		0.0008452 \times 1000 = 0.845
	\end{align*}
	On average, a \$1,000 increase in income is associated with 0.845 unit increase in prestige score, for professional jobs.
	
	
	
	\item [(g)]
	What is the effect of changing one's occupations from non-professional to professional when her income is \$6,000? We are interested in the marginal effect of professional jobs when the variable \texttt{income} takes the value of $6,000$. Calculate the change in $\hat{y}$ based on your answer for (c).\\
	
	Answer: \\
	General prediction equation.
	\begin{align*}
		\hat{prestige} &= 21.1423 + 0.003171income + 37.7813professional - 0.002326income \times professional
	\end{align*}

	\lstinputlisting[language=R, firstline=48, lastline=51]{PS04_answersJX.R}
	On average, when income is \$6,000, changing from a non-professional job to a professional job is associated with 23.83 unit increase in prestige score.
	
	
	\end{enumerate}

\newpage

\section*{Question 2: Political Science}
\vspace{.25cm}
\noindent 	Researchers are interested in learning the effect of all of those yard signs on voting preferences.\footnote{Donald P. Green, Jonathan	S. Krasno, Alexander Coppock, Benjamin D. Farrer,	Brandon Lenoir, Joshua N. Zingher. 2016. ``The effects of lawn signs on vote outcomes: Results from four randomized field experiments.'' Electoral Studies 41: 143-150. } Working with a campaign in Fairfax County, Virginia, 131 precincts were randomly divided into a treatment and control group. In 30 precincts, signs were posted around the precinct that read, ``For Sale: Terry McAuliffe. Don't Sellout Virgina on November 5.'' \\

Below is the result of a regression with two variables and a constant.  The dependent variable is the proportion of the vote that went to McAuliff's opponent Ken Cuccinelli. The first variable indicates whether a precinct was randomly assigned to have the sign against McAuliffe posted. The second variable indicates
a precinct that was adjacent to a precinct in the treatment group (since people in those precincts might be exposed to the signs).  \\

\vspace{.5cm}
\begin{table}[!htbp]
	\centering 
	\textbf{Impact of lawn signs on vote share}\\
	\begin{tabular}{@{\extracolsep{5pt}}lccc} 
		\\[-1.8ex] 
		\hline \\[-1.8ex]
		Precinct assigned lawn signs  (n=30)  & 0.042\\
		& (0.016) \\
		Precinct adjacent to lawn signs (n=76) & 0.042 \\
		&  (0.013) \\
		Constant  & 0.302\\
		& (0.011)
		\\
		\hline \\
	\end{tabular}\\
	\footnotesize{\textit{Notes:} $R^2$=0.094, N=131}
\end{table}

\vspace{.5cm}
\begin{enumerate}
	\item [(a)] Use the results from a linear regression to determine whether having these yard signs in a precinct affects vote share (e.g., conduct a hypothesis test with $\alpha = .05$).\\
	\newpage
	Answer: \\
	Null hypothesis: Having these yard signs in a precinct has 0 effect to vote share.\\
	Alternative: Having these yard signs in a precinct has some effects to vote share.\\
	\begin{align*}
		H_0: \beta_1 = 0 \quad \text{vs.} \quad H_a: \beta_1 \neq 0
	\end{align*}
	
	Test statistic:
	\begin{align*}
		t = \frac{\hat{\beta}_1 - 0}{se_{\hat{\beta}_1}} = \frac{0.042 - 0}{0.016} = 2.625
	\end{align*}
	
	p-value:
	\begin{align}
		df = n-k-1 = 131-2-1 = 128
	\end{align}
	\lstinputlisting[language=R, firstline=54, lastline=56]{PS04_answersJX.R}
	\begin{align}
		p = 0.00972002, \quad p < \alpha
	\end{align}
	
	Conclusion: Reject null, there is enough evidence to conclude that there having these yard signs in a precinct has some effects on vote share.
	
	
	\newpage		
	\item [(b)]  Use the results to determine whether being
	next to precincts with these yard signs affects vote
	share (e.g., conduct a hypothesis test with $\alpha = .05$). \\
	
	Answer: \\
	Null hypothesis: Being next to precincts with these yard signs have 0 effect on vote share.\\
	Alternative: Being next to precincts with these yard signs have some effects on vote share.\\
	\begin{align*}
		H_0: \beta_2 = 0 \quad \text{vs.} \quad H_a: \beta_2 \neq 0
	\end{align*}
	
	Test statistic:
	\begin{align*}
		t = \frac{\hat{\beta}_2 - 0}{se_{\hat{\beta}_2}} = \frac{0.042 - 0}{0.013} = 3.2308
	\end{align*}
	
	p-value:
	\begin{align}
		df = n-k-1 = 131-2-1 = 128
	\end{align}
	\lstinputlisting[language=R, firstline=59, lastline=61]{PS04_answersJX.R}
	\begin{align}
		p = 0.00156946, \quad p < \alpha
	\end{align}
	
	Conclusion: Reject null, there is enough evidence to conclude that being next to precincts with these yard signs have some effects to vote share.
	
	\newpage
	\item [(c)] Interpret the coefficient for the constant term substantively.\\
	\begin{doublespace}
		Answers: \\
		For precincts that were neither assigned lawn signs nor adjacent to a precinct with lawn signs, the predicted vote share for Ken Cuccinelli is 0.302, or about 30.2\%, on average.
	\end{doublespace}
	
	\item [(d)] Evaluate the model fit for this regression.  What does this	tell us about the importance of yard signs versus other factors that are not modeled? \\
	\begin{doublespace}
		Answers: \\
		The $R^2$ of the model is 0.094, it is low and meaning that the signs only explains a small port of why precincts differ in vote share. This tells us the yard signs are less important compare with partyID, income or other factors that are not modeled. Although $R^2$ is low, the treatment is randomly assigned, the significant coefficient still indicate that the yard sign has an effect.
	\end{doublespace}
	
\end{enumerate}  


\end{document}
